% GENERATED FILE. Edit canonical lessons, then run npm run generate:books.
% canonical-source-hash: fnv1a64:b6b4ffc0122bab97
% canonical-lessons: BN-C366-sombhabona, BN-C366-sombhoboto, BN-C366-onibarjo, BN-C366-birol, BN-C366-dhore-neoya

\chapter{Likely or Not}
\label{ch:bn-likely-or-not}
\hlchaptermodality{366}

\begin{chapteropening}
\textbf{By the end of this chapter:} \emph{I can say how likely something is and what follows if it happens.}

Complete the last lesson of chapter 366: I can say how likely something is and what follows if it happens.
\end{chapteropening}

\section[\texorpdfstring{\bn{সম্ভাবনা} (sômbhābônā)}{সম্ভাবনা (sômbhābônā)}]{\bn{সম্ভাবনা} (sômbhābônā) — a likelihood; a chance}

\label{lesson:BN-C366-sombhabona}



\begin{quote}
Before the new one: say the Bengali for to talk round, then the Bengali for unanimous.
\end{quote}

\subsection*{You'll want to know: \bn{সম্ভাবনা}}
\textbf{\bn{সম্ভাবনা}} — \emph{sômbhābônā} — \textquotedblleft{}a likelihood; a chance\textquotedblright{}.

\begin{quote}
\bn{আজ} \bn{ঝড়ের} \bn{সম্ভাবনা} \bn{আছে}, \bn{তাই} \bn{বাইরে} \bn{না} \bn{গেলে} \bn{ভালো।}
\end{quote}

\begin{quote}
\emph{āj jhôṛer sômbhābônā āchhe, tāi bāire nā gele bhālo.}
\end{quote}

\begin{quote}
There's a chance of a storm today, so it's better not to go out.
\end{quote}

\begin{cousinweb}
\textbf{\bn{সম্ভাবনা}} (\emph{sômbhābônā}) is Sanskrit \emph{sambhāvanā}, supposing or assuming (Monier-Williams), from the same root as \bn{সম্ভব}, possible. Bengali uses it for how likely a thing is.
\end{cousinweb}

\begin{grammarlens}[title={What you've built}]
One more word for saying what happens if.
\end{grammarlens}

\subsection*{Your turn}
\begin{itemize}
\raggedright
  \item \emph{Say it:} \emph{sômbhābônā}
  \item \emph{Say it:} \emph{sômbhābônā}, once more
  \item \emph{Say it:} \emph{āj jhôṛer sômbhābônā āchhe, tāi bāire nā gele bhālo.} — There's a chance of a storm today, so it's better not to go out.
  \item \emph{Recall:} say the Bengali for to talk round, then the Bengali for unanimous, then say \emph{sômbhābônā} again
\end{itemize}

\begin{tcolorbox}[breakable,colback=teal!4,colframe=teal!35!black,title={Before you move on}]
What does \bn{সম্ভাবনা} mean? (\textquotedblleft{}A likelihood; a chance\textquotedblright{}.) Say it once more.
\end{tcolorbox}
\section[\texorpdfstring{\bn{সম্ভবত} (sômbhôbôtô)}{সম্ভবত (sômbhôbôtô)}]{\bn{সম্ভবত} (sômbhôbôtô) — probably}

\label{lesson:BN-C366-sombhoboto}



\begin{quote}
Before the new one: say the Bengali for unanimous, then the Bengali for a likelihood.
\end{quote}

\subsection*{You'll want to know: \bn{সম্ভবত}}
\textbf{\bn{সম্ভবত}} — \emph{sômbhôbôtô} — \textquotedblleft{}probably\textquotedblright{}.

\begin{quote}
\bn{অরুণ} \bn{যদি} \bn{বাস} \bn{না} \bn{পায়}, \bn{তাহলে} \bn{সম্ভবত} \bn{দেরি} \bn{করবে।}
\end{quote}

\begin{quote}
\emph{ôrun jôdi bās nā pāy, tāhôle sômbhôbôtô deri kôrbe.}
\end{quote}

\begin{quote}
If Arun doesn't get a bus, he will probably be late.
\end{quote}

\begin{cousinweb}
\textbf{\bn{সম্ভবত}} (\emph{sômbhôbôtô}) is Sanskrit \emph{sambhava}, arising or possible (Monier-Williams), with the adverb ending \emph{-tas}, in its Bengali form: 'in all likelihood'. It sits beside the taught \bn{সম্ভব}.
\end{cousinweb}

\begin{grammarlens}[title={What you've built}]
One more word for saying what happens if.
\end{grammarlens}

\subsection*{Your turn}
\begin{itemize}
\raggedright
  \item \emph{Say it:} \emph{sômbhôbôtô}
  \item \emph{Say it:} \emph{sômbhôbôtô}, once more
  \item \emph{Say it:} \emph{ôrun jôdi bās nā pāy, tāhôle sômbhôbôtô deri kôrbe.} — If Arun doesn't get a bus, he will probably be late.
  \item \emph{Recall:} say the Bengali for unanimous, then the Bengali for a likelihood, then say \emph{sômbhôbôtô} again
\end{itemize}

\begin{tcolorbox}[breakable,colback=teal!4,colframe=teal!35!black,title={Before you move on}]
What does \bn{সম্ভবত} mean? (\textquotedblleft{}Probably\textquotedblright{}.) Say it once more.
\end{tcolorbox}
\section[\texorpdfstring{\bn{অনিবার্য} (ônibārjô)}{অনিবার্য (ônibārjô)}]{\bn{অনিবার্য} (ônibārjô) — inevitable; unavoidable}

\label{lesson:BN-C366-onibarjo}



\begin{quote}
Before the new one: say the Bengali for a likelihood, then the Bengali for probably.
\end{quote}

\subsection*{You'll want to know: \bn{অনিবার্য}}
\textbf{\bn{অনিবার্য}} — \emph{ônibārjô} — \textquotedblleft{}inevitable; unavoidable\textquotedblright{}.

\begin{quote}
\bn{নিয়ম} \bn{না} \bn{মানলে} \bn{দুর্ঘটনা} \bn{অনিবার্য।}
\end{quote}

\begin{quote}
\emph{niyôm nā mānle durghôṭônā ônibārjô.}
\end{quote}

\begin{quote}
If people don't follow the rules, an accident is inevitable.
\end{quote}

\begin{cousinweb}
\textbf{\bn{অনিবার্য}} (\emph{ônibārjô}) is Sanskrit \emph{anivārya}, not to be warded off, unavoidable (Monier-Williams), with the negative \emph{a-} at its head: what cannot be held back. It marks the likeliest outcome of all, the one that will surely follow.
\end{cousinweb}

\begin{grammarlens}[title={What you've built}]
One more word for saying what happens if.
\end{grammarlens}

\subsection*{Your turn}
\begin{itemize}
\raggedright
  \item \emph{Say it:} \emph{ônibārjô}
  \item \emph{Say it:} \emph{ônibārjô}, once more
  \item \emph{Say it:} \emph{niyôm nā mānle durghôṭônā ônibārjô.} — If people don't follow the rules, an accident is inevitable.
  \item \emph{Recall:} say the Bengali for a likelihood, then the Bengali for probably, then say \emph{ônibārjô} again
\end{itemize}

\begin{tcolorbox}[breakable,colback=teal!4,colframe=teal!35!black,title={Before you move on}]
What does \bn{অনিবার্য} mean? (\textquotedblleft{}Inevitable; unavoidable\textquotedblright{}.) Say it once more.
\end{tcolorbox}
\section[\texorpdfstring{\bn{বিরল} (birôl)}{বিরল (birôl)}]{\bn{বিরল} (birôl) — rare}

\label{lesson:BN-C366-birol}



\begin{quote}
Before the new one: say the Bengali for probably, then the Bengali for inevitable.
\end{quote}

\subsection*{You'll want to know: \bn{বিরল}}
\textbf{\bn{বিরল}} — \emph{birôl} — \textquotedblleft{}rare\textquotedblright{}.

\begin{quote}
\bn{এই} \bn{পাখিটা} \bn{বিরল}, \bn{দেখতে} \bn{পেলে} \bn{ছবি} \bn{তুলবে।}
\end{quote}

\begin{quote}
\emph{ei pākhiṭā birôl, dekhte pele chhôbi tulbe.}
\end{quote}

\begin{quote}
This bird is rare, so if you get to see it, take a photo.
\end{quote}

\begin{cousinweb}
\textbf{\bn{বিরল}} (\emph{birôl}) is Sanskrit \emph{virala}, rare or thinly scattered (Monier-Williams).
\end{cousinweb}

\begin{grammarlens}[title={What you've built}]
One more word for saying what happens if.
\end{grammarlens}

\subsection*{Your turn}
\begin{itemize}
\raggedright
  \item \emph{Say it:} \emph{birôl}
  \item \emph{Say it:} \emph{birôl}, once more
  \item \emph{Say it:} \emph{ei pākhiṭā birôl, dekhte pele chhôbi tulbe.} — This bird is rare, so if you get to see it, take a photo.
  \item \emph{Recall:} say the Bengali for probably, then the Bengali for inevitable, then say \emph{birôl} again
\end{itemize}

\begin{tcolorbox}[breakable,colback=teal!4,colframe=teal!35!black,title={Before you move on}]
What does \bn{বিরল} mean? (\textquotedblleft{}Rare\textquotedblright{}.) Say it once more.
\end{tcolorbox}
\section[\texorpdfstring{\bn{ধরে} \bn{নেওয়া} (dhôre neoyā)}{ধরে নেওয়া (dhôre neoyā)}]{\bn{ধরে} \bn{নেওয়া} (dhôre neoyā) — to assume; to take for granted}

\label{lesson:BN-C366-dhore-neoya}



\begin{quote}
Before the new one: say the Bengali for inevitable, then the Bengali for rare.
\end{quote}

\subsection*{You'll want to know: \bn{ধরে} \bn{নেওয়া}}
\textbf{\bn{ধরে} \bn{নেওয়া}} — \emph{dhôre neoyā} — \textquotedblleft{}to assume; to take for granted\textquotedblright{}.

\begin{quote}
\bn{যদি} \bn{ধরে} \bn{নিই} \bn{যে} \bn{কাল} \bn{ছুটি}, \bn{তাহলে} \bn{আমরা} \bn{কোথায়} \bn{যাব}?
\end{quote}

\begin{quote}
\emph{jôdi dhôre nii je kāl chhuṭi, tāhôle āmrā kothāy jābô?}
\end{quote}

\begin{quote}
If we assume tomorrow is a holiday, where shall we go?
\end{quote}

Past (he/she …): \bn{তিনি} \bn{ধরে} \bn{নিলেন} \bn{যে} \bn{সবাই} \bn{আসবে}, \emph{tini dhôre nilen je sôbāi āsbe}.

\begin{cousinweb}
\textbf{\bn{ধরে} \bn{নেওয়া}} (\emph{dhôre neoyā}) is \bn{ধরে}, the -\bn{এ} form of \bn{ধরা} 'to hold' (CDIAL 6747), with \bn{নেওয়া} 'to take' (CDIAL 6966): to take a thing as held, to assume it. What is assumed follows \bn{যে}.
\end{cousinweb}

\begin{grammarlens}[title={What you've built}]
One more word for saying what happens if.
\end{grammarlens}

\subsection*{Your turn}
\begin{itemize}
\raggedright
  \item \emph{Say it:} \emph{dhôre neoyā}
  \item \emph{Say it:} \emph{dhôre neoyā}, once more
  \item \emph{Say it:} \emph{jôdi dhôre nii je kāl chhuṭi, tāhôle āmrā kothāy jābô?} — If we assume tomorrow is a holiday, where shall we go?
  \item \emph{Recall:} say the Bengali for inevitable, then the Bengali for rare, then say \emph{dhôre neoyā} again
\end{itemize}

\begin{tcolorbox}[breakable,colback=teal!4,colframe=teal!35!black,title={Before you move on}]
What does \bn{ধরে} \bn{নেওয়া} mean? (\textquotedblleft{}To assume; to take for granted\textquotedblright{}.) Say it once more.
\end{tcolorbox}
